\documentclass{article}

\usepackage{amsmath,amssymb}
\usepackage{xurl}
\usepackage[hidelinks]{hyperref}
\setlength{\emergencystretch}{2em}

% Shared commands for notes in docs/paper-gaps/.

\DeclareUnicodeCharacter{2080}{\ifmmode{}_{0}\else\textsubscript{0}\fi}
\DeclareUnicodeCharacter{2081}{\ifmmode{}_{1}\else\textsubscript{1}\fi}
\DeclareUnicodeCharacter{2082}{\ifmmode{}_{2}\else\textsubscript{2}\fi}
\DeclareUnicodeCharacter{2099}{\ifmmode{}_{n}\else\textsubscript{n}\fi}

\newcommand{\Lean}{\textsc{Lean}}
\ExplSyntaxOn
\NewDocumentCommand{\leanid}{m}{
  \ifmmode
    \textsf{\detokenize{#1}}
  \else
    \group_begin:
    \urlstyle{sf}
    \tl_set:Nn \l_tmpa_tl {#1}
    \tl_replace_all:Nnn \l_tmpa_tl {\_} {_}
    \exp_args:NV \path \l_tmpa_tl
    \group_end:
  \fi
}
\ExplSyntaxOff
\providecommand{\tr}{\operatorname{tr}}
\newcommand{\tnleanIssueBase}{https://github.com/LionSR/TNLean/issues/}
\newcommand{\tnleanPrBase}{https://github.com/LionSR/TNLean/pull/}
\newcommand{\tnleanDocBase}{https://sirui-lu.com/TNLean/docs/}
\newcommand{\ghissue}[1]{\href{\tnleanIssueBase#1}{GitHub issue \##1}}
\newcommand{\ghpr}[1]{\href{\tnleanPrBase#1}{PR \##1}}
\newcommand{\doclink}[1]{\href{\tnleanDocBase#1.html}{\path{#1}}}

% Dirac notation and the identity operator, as in blueprint/src/macros/common.tex.
% \providecommand keeps note-local definitions authoritative.
\usepackage{dsfont}
\providecommand{\ket}[1]{|#1\rangle}
\providecommand{\bra}[1]{\langle #1|}
\providecommand{\braket}[2]{\langle #1 | #2 \rangle}
\providecommand{\ketbra}[2]{|#1\rangle\!\langle #2|}
\providecommand{\Id}{\mathds{1}}
\providecommand{\one}{\mathds{1}}
\providecommand{\C}{\mathbb{C}}
\providecommand{\MN}[1]{M_{#1}(\C)}
\providecommand{\MD}{M_D(\C)}

% External referencing for paper sources.  The build helper writes sanitized
% auxiliary files under build/paper-aux/ and excludes duplicated source labels.
\usepackage{xr}

% Import a generated paper auxiliary file with a caller-supplied label prefix.
% The candidate paths support builds from docs/paper-gaps/, the repository root,
% and build/paper-aux/ itself.
\newcommand{\importpaperaux}[2]{%
  \IfFileExists{../../build/paper-aux/#2.aux}{%
    \externaldocument[#1]{../../build/paper-aux/#2}%
  }{%
    \IfFileExists{build/paper-aux/#2.aux}{%
      \externaldocument[#1]{build/paper-aux/#2}%
    }{%
      \IfFileExists{#2.aux}{%
        \externaldocument[#1]{#2}%
      }{}%
    }%
  }%
}

% General external reference macros
\newcommand{\paperref}[2]{\ref{#1-#2}}
\newcommand{\papereqref}[2]{(\ref{#1-#2})}

% CPSV16 source (1606.00608 / MPDO-22-12-17-2.tex) external document and shortcuts
\importpaperaux{cpsv16-}{1606.00608/MPDO-22-12-17-2-tnlean}
\newcommand{\cpsvref}[1]{\paperref{cpsv16}{#1}}
\newcommand{\cpsveqref}[1]{\papereqref{cpsv16}{#1}}

% Verdict marker: \gapnote{<kind>}{<status>}. Kind is the policy
% classification (clarification, local-correction, scope-restriction,
% unfaithful, false-source, open-gap); status is open, wip (elimination
% underway), resolved, or historical. Machine-read by the site index;
% prints a verdict line.
\newcommand{\gapnote}[2]{\par\noindent\textbf{Verdict:} #1 (#2).\par}


\title{Normal Parent-Hamiltonian Range Reduction in Cirac--Perez-Garcia--Schuch--Verstraete 2021}
\author{}
\date{2026-04-30}

\begin{document}
\maketitle
\gapnote{scope-restriction}{resolved}

\section{The assertion}

Section~IV.C of Cirac--P\'erez-Garc\'ia--Schuch--Verstraete proves a
range-reduction result for normal matrix product states
\cite{Cirac2021Matrix}.\footnote{For traceability, this note corresponds to
\href{https://github.com/LionSR/TNLean/issues/1042}{issue \#1042}. The principal
range-reduction comparison is recorded in
\href{https://github.com/LionSR/TNLean/issues/588}{issue \#588}; the
periodic-boundary comparison is recorded in
\href{https://github.com/LionSR/TNLean/issues/2368}{issue \#2368}; its current
coordinate form was tracked by
\href{https://github.com/LionSR/TNLean/issues/2405}{issue \#2405}, now closed
after the length-$L_0$ trace-probe step isolated in
\href{https://github.com/LionSR/TNLean/issues/2929}{issue \#2929}; and the
parent-Hamiltonian consequences are recorded in
\href{https://github.com/LionSR/TNLean/issues/460}{issue \#460}.}

Let $A$ be a normal MPS tensor, and let $L_0>0$ be an injectivity length for
$A$: blocking $L_0$ consecutive sites gives an injective tensor. For a
contiguous interval $I$, write $\mathcal G_I(A)$ for its local MPS ground
space. The source considers terms of range $L_0+1$ supported on sites
$1,\ldots,L_0+1$ and $2,\ldots,L_0+2$. It partitions their union into the
first site, the middle sites $2,\ldots,L_0+1$, and the last site. The middle
region has length $L_0$, so its blocked tensor is injective.

The source inverts the tensor on the first-site region, restores the middle
region, and uses injectivity on the enlarged region to prove
\begin{align}
    \mathcal G_{\{1,\ldots,L_0+1\}}(A)
        \cap \mathcal G_{\{2,\ldots,L_0+2\}}(A)
        &= \mathcal G_{\{1,\ldots,L_0+2\}}(A).
    \label{eq:cpgsv21_normal_range_reduction_intersection_identity}
\end{align}
The reverse inclusion in
Eq.~\ref{eq:cpgsv21_normal_range_reduction_intersection_identity} is
immediate.\footnote{Source location:
\path{Papers/2011.12127/TN-Review-main.tex}:2049--2077, especially the
middle-region argument around the displayed figures following
\path{eq:4:initiate-intersection-proof}.}

Iterating this identity shows that the ground space defined by all terms of
length $L_0+1$ agrees with the ground space defined by all terms of length
$L_0+k$. For $k=L_0$, the source invokes either the preceding range-$2L_0$
uniqueness theorem or the analogous periodic-boundary argument. It concludes
that a normal MPS which becomes injective after blocking $L_0$ sites is the
unique ground state of its parent Hamiltonian with interaction range $L_0+1$
\cite{Cirac2021Matrix}.\footnote{Source location:
\path{Papers/2011.12127/TN-Review-main.tex}:2078--2079 for the iteration and
periodic-boundary closure-property sentence, and
\path{Papers/2011.12127/TN-Review-main.tex}:2087--2090 for the theorem labelled
\path{thm:4:unique-gs-L0_plus_1}.}

\section{The formal statement}

For a ring of length $N$ and interaction range $L$, write
$\mathcal G_{N,L}(A)$ for the periodic-boundary ground space, and write
$\Omega_N(A)$ for the periodic MPS vector. The formal theorem assumes that
$A$ is normal and injective after blocking $L_0>0$ sites, with
\begin{align}
    2 &\leq N,
    \label{eq:cpgsv21_normal_range_reduction_ring_lower_bound}\\
    L_0+1 &\leq N,
    \label{eq:cpgsv21_normal_range_reduction_source_endpoint}\\
    L_0 &< L\leq N.
    \label{eq:cpgsv21_normal_range_reduction_interaction_range}
\end{align}
It proves
\begin{align}
    \mathcal G_{N,L}(A)
        &= \mathbb C\,\Omega_N(A).
    \label{eq:cpgsv21_normal_range_reduction_periodic_ground_space}
\end{align}
Thus the formal statement retains the source endpoint $L_0+1\leq N$. The
difficult inclusion in
Eq.~\ref{eq:cpgsv21_normal_range_reduction_periodic_ground_space} combines a
complementary-word argument for $L_0+1<N$ with a full-ring cyclic change of cut
for $N=L_0+1$.\footnote{Formal locations:
\path{TNLean/MPS/ParentHamiltonian/UniqueGroundState.lean}:738 for
\leanid{MPSTensor.chainGroundSpace_le_mpvSubmodule_of_normal_range_reduction},
and
\path{TNLean/MPS/ParentHamiltonian/UniqueGroundState.lean}:776 for
\leanid{MPSTensor.chainGroundSpace_eq_mpvSubmodule_normal}. The formal
hypotheses are named \leanid{Kraus.IsNormal A} and
\leanid{Kraus.IsNBlkInjective A L0}, together with
\path{0 < L0}, \path{2 <= N}, \path{L0 + 1 <= N}, and
\path{L0 < L <= N}.}

The blueprint records the closure-property step in
\path{blueprint/src/chapter/ch13_parent_hamiltonian_normal_closing_reverse.tex}
under the label \path{thm:normal_range_reduction_containment}. Its proof sketch
separates the open-chain iteration from the periodic-boundary comparison, as
does the source argument \cite{Cirac2021Matrix}.

The first formal mismatch concerns the proof through the middle region. The
available intersection theorem applies to a tensor that is already injective
at one site. Its proof writes the identity as a linear combination of the
one-site matrices $A_j$ and iterates the resulting single-letter intersection
property through contiguous windows.\footnote{Formal locations:
\path{TNLean/MPS/ParentHamiltonian/IntersectionProperty.lean}:343 for
\leanid{MPSTensor.groundSpace_intersection},
\path{TNLean/MPS/Chain/OneSidedInverse.lean}:67 for
\leanid{Kraus.decompositionMap hA 1}, and
\path{TNLean/MPS/ParentHamiltonian/CyclicWindow.lean}:220--257 for the
contiguous iteration.}
For a normal tensor with injectivity length $L_0$, the available inverse acts
on words of length $L_0$, not on individual letters. This does not contradict
the source; it shows that the diagrammatic region-by-region proof is not a
literal instance of the older formal intersection theorem.

The formal argument treats the open and periodic parts separately. A
block-injective range-reduction lemma proves that if every contiguous window
of length $L_0+1$ belongs to its local MPS ground space, then the whole open
chain belongs to its full MPS ground space. Cyclic-window monotonicity derives
this conclusion from the cyclic length-$L$ constraints. This formalizes the
iterated open-boundary argument without using the single-site intersection
lemma.\footnote{Formal locations:
\path{TNLean/MPS/ParentHamiltonian/UniqueGroundState.lean}:107 for
\leanid{MPSTensor.contiguous_mem_groundSpace_of_isNBlkInjective}, and
\path{TNLean/MPS/ParentHamiltonian/UniqueGroundState.lean}:169 for
\leanid{MPSTensor.chainGroundSpace_le_groundSpace_of_isNBlkInjective}.}

\section{The boundary block-window identity}

For a word $\omega$, let $A^\omega$ denote the ordered product of its tensor
matrices. After open-boundary range reduction, a periodic-boundary ground
state has a trace-form representation with a matrix $X$. The boundary
comparison must provide, for every word $\alpha$ of length $L_0$ and every
nonempty complementary word $\nu$, a matrix $Y_\nu$ satisfying
\begin{align}
    XA^\alpha A^\nu
        &= A^\alpha Y_\nu.
    \label{eq:cpgsv21_normal_range_reduction_block_window_matrix_identity}
\end{align}
The block-injective commutation lemma then gives
\begin{align}
    XA^k
        &= A^kX
        \qquad\text{for every physical letter $k$}.
    \label{eq:cpgsv21_normal_range_reduction_one_site_commutation}
\end{align}
Block injectivity finally forces $X$ to be scalar, so the state belongs to
$\mathbb C\,\Omega_N(A)$.

Define the length-$N$ trace-form vector $\Gamma_N(X)$ by
\begin{align}
    \Gamma_N(X)_\omega
        &= \tr(A^\omega X).
    \label{eq:cpgsv21_normal_range_reduction_trace_form_vector}
\end{align}
Membership of $\Gamma_N(X)$ in the open-chain space does not imply membership
in the periodic-boundary ground space. For a boundary-crossing window, the
restricted coefficients have the form
\begin{align}
    \tr\!\left(A^{\omega_{\mathrm{tail}}}X
        A^{\omega_{\mathrm{head}}}A^\rho\right).
    \label{eq:cpgsv21_normal_range_reduction_boundary_crossing_coefficient}
\end{align}
Membership in the length-$L$ local ground space instead requires matrices
$Y_\rho$, depending only on labels outside the window, such that these
coefficients equal
\begin{align}
    \tr\!\left(A^{\omega_{\mathrm{tail}}}
        A^{\omega_{\mathrm{head}}}Y_\rho\right).
    \label{eq:cpgsv21_normal_range_reduction_local_ground_space_coefficient}
\end{align}
The periodic closure-property comparison moves $X$ through the
boundary-crossing word and supplies this equality.

In the coordinate proof, $\mu$ is the nonempty word outside the two cyclic
windows, and $\eta$ is the letter on the remaining outside sites. This realizes
the closure-property statement in Section~IV.C, lines~2078--2079, of
Ref.~\cite{Cirac2021Matrix}. Choose boundary-restriction matrices $W_\eta$ and
$V_\eta$ for the two boundary-crossing restrictions. The one-sided
boundary-product lemma gives
\begin{align}
    W_\eta A^j
        &= A^\mu A^jX,
    \label{eq:cpgsv21_normal_range_reduction_left_boundary_product}\\
    XA^jA^\mu
        &= A^jV_\eta
        \qquad\text{for every physical letter $j$}.
    \label{eq:cpgsv21_normal_range_reduction_right_boundary_product}
\end{align}
Here $W_\eta$ corresponds to the restriction beginning at the last site, and
$V_\eta$ corresponds to the second boundary-crossing restriction.\footnote{
Formal locations:
\path{TNLean/MPS/ParentHamiltonian/BoundaryClosingStripping.lean} for
\leanid{MPSTensor.closure_property_boundary_restrictions_eq_of_groundSpaceMap},
and
\path{TNLean/MPS/ParentHamiltonian/BoundaryClosing.lean} for
\leanid{MPSTensor.closure_property_boundary_one_sided_products_of_groundSpaceMap}.}
The block-injective commutation lemma applied to this local comparison yields
Eq.~\ref{eq:cpgsv21_normal_range_reduction_one_site_commutation}.\footnote{In
the proof of
\leanid{MPSTensor.closure_property_boundary_restrictions_eq_of_groundSpaceMap},
the commutation identity follows from
\leanid{MPSTensor.boundary_matrix_commutes_of_isNBlkInjective_of_block_matEq}.}

For an $N$-site vector $\psi$, let
$\operatorname{Res}^{\tau}_{i,L}(\psi)$ denote the $L$-site vector whose
$\sigma$-coordinate is obtained by placing $\sigma$ in the cyclic window of
length $L$ beginning at $i$ and placing $\tau$ on the remaining sites. Once
Eq.~\ref{eq:cpgsv21_normal_range_reduction_one_site_commutation} is available,
the boundary-restriction equality lemma proves
\begin{align}
    W_\eta A^j
        &= V_\eta A^j
        \qquad\text{for every $\eta$ and $j$}.
    \label{eq:cpgsv21_normal_range_reduction_boundary_first_letter_products}
\end{align}
The first-letter restriction lemma then gives
\begin{align}
    \operatorname{Res}^{\tau^+_\eta(\mu)}_{M,L_0+1}(\psi)
        &=
    \operatorname{Res}^{\tau^-_\eta(\mu)}_{M+1-L_0,L_0+1}(\psi).
    \label{eq:cpgsv21_normal_range_reduction_boundary_restriction_equality}
\end{align}
Thus Eq.~\ref{eq:cpgsv21_normal_range_reduction_boundary_restriction_equality}
follows from the one-sided products and the commutation identity; it is not
the primitive comparison.\footnote{Formal locations:
\path{TNLean/MPS/ParentHamiltonian/BoundaryMatrixBlock.lean} for
\leanid{MPSTensor.boundary_restrictions_eq_of_commutes_and_one_sided}, and
\path{TNLean/MPS/ParentHamiltonian/BoundaryClosingStripping.lean} for its use in
\leanid{MPSTensor.closure_property_boundary_restrictions_eq_of_groundSpaceMap}.}

The block-injective commutation lemma reduces
Eq.~\ref{eq:cpgsv21_normal_range_reduction_one_site_commutation} to
Eq.~\ref{eq:cpgsv21_normal_range_reduction_block_window_matrix_identity}.
\footnote{Formal location:
\path{TNLean/MPS/ParentHamiltonian/BoundaryMatrixBlock.lean} for
\leanid{MPSTensor.boundary_matrix_commutes_of_isNBlkInjective_of_block_matEq}.}
The last cyclic window first yields only
\begin{align}
    \tr\!\left(A^jXA^\alpha A^\nu\right)
        &=
    \tr\!\left(A^jA^\alpha Y_\nu\right).
    \label{eq:cpgsv21_normal_range_reduction_single_letter_trace_rotation}
\end{align}
This identity is recorded as
\leanid{MPSTensor.closure_property_boundary_block_window_trace_eq_of_groundSpaceMap}.
Single-letter probes do not separate matrices for a general
$L_0$-block-injective tensor. The same cyclic-window comparison therefore
establishes, for all length-$L_0$ words $\alpha,\beta$ and every nonempty
complementary word $\nu$,
\begin{align}
    \tr\!\left(A^\beta XA^\alpha A^\nu\right)
        &=
    \tr\!\left(A^\beta A^\alpha Y_\nu\right).
    \label{eq:cpgsv21_normal_range_reduction_block_word_trace_identity}
\end{align}
These identities are formalized by
\leanid{MPSTensor.closure_property_boundary_block_window_trace_evalWord_mul_eq_of_groundSpaceMap}.
Because the products $A^\beta$ span the full matrix algebra, trace separation
turns Eq.~\ref{eq:cpgsv21_normal_range_reduction_block_word_trace_identity}
into
Eq.~\ref{eq:cpgsv21_normal_range_reduction_block_window_matrix_identity}.
\footnote{Formal location:
\path{TNLean/MPS/ParentHamiltonian/BoundaryClosingStripping.lean} for
\leanid{MPSTensor.block_window_matrix_equation_of_trace_evalWord_mul_eq_of_isNBlkInjective}.}
The source-labelled formal statement is
\leanid{MPSTensor.closure_property_boundary_block_window_equation_of_groundSpaceMap}.

The same comparison has a coordinate form:
\begin{align}
    \tau^{+}_{\eta}(\mu)_i
        &=
        \begin{cases}
            \mu_{i-L_0}, & L_0\leq i<N-1,\\
            \eta, & \text{otherwise},
        \end{cases}
    \label{eq:cpgsv21_normal_range_reduction_positive_boundary_word}\\
    \tau^{-}_{\eta}(\mu)_i
        &=
        \begin{cases}
            \mu_{i-1}, & 1\leq i<N-L_0,\\
            \eta, & \text{otherwise}.
        \end{cases}
    \label{eq:cpgsv21_normal_range_reduction_negative_boundary_word}
\end{align}
The direct comparison of the corresponding restrictions, and the
left-multiplied comparison obtained from first-letter restrictions, follow
from Eq.~\ref{eq:cpgsv21_normal_range_reduction_one_site_commutation}.

\section{Resolved endpoint: the minimal ring}
\label{sec:cpgsv21_normal_range_reduction_minimal_ring}

The source theorem includes $N=L_0+1$ \cite{Cirac2021Matrix}. This case has no
nonempty complementary word and is proved separately. Write a vector in the
full-ring ground space as $\Gamma_{L_0+1}(X)$. Moving the cut by $L_0$ sites
and expressing the translated state as $\Gamma_{L_0+1}(Y)$ gives
\begin{align}
    XA^j
        &= A^jY
        \qquad\text{for every physical letter $j$}.
    \label{eq:cpgsv21_normal_range_reduction_full_ring_intertwining}
\end{align}
Repeating this comparison around all $L_0+1$ cuts shows that $X$ commutes with
every full-ring word. Block injectivity persists at length $L_0+1$, so these
words span the full matrix algebra. Hence $X$ is scalar and
\begin{align}
    \mathcal G_{L_0+1,L_0+1}(A)
        &= \mathbb C\,\Omega_{L_0+1}(A).
    \label{eq:cpgsv21_normal_range_reduction_minimal_ring_ground_space}
\end{align}
This result also removes the extra strict bound from the range-$2L_0$
corollary. In particular, it covers $L_0=1$ and $N=2$ directly.

The formal proof applies
\leanid{MPSTensor.full_ring_boundary_intertwines_of_chainGroundSpace} and then
\leanid{MPSTensor.chainGroundSpace_le_mpvSubmodule_of_isNBlkInjective_full_ring}.
The public theorems
\leanid{MPSTensor.chainGroundSpace_le_mpvSubmodule_of_normal_range_reduction},
\leanid{MPSTensor.chainGroundSpace_eq_mpvSubmodule_normal},
\leanid{MPSTensor.parentHamiltonian_unique_gs_normal}, and
\leanid{MPSTensor.parentHamiltonian_unique_gs_injective} now use the
non-strict source bound.

\section{Conclusion}

The source argument is not refuted. Ref.~\cite{Cirac2021Matrix} compresses two
steps that the formal proof separates: the open-chain range reduction obtained
by inverting and restoring the middle region, and the movement of the
trace-form matrix $X$ through a boundary-crossing word.

For $L_0+1<N$, the cyclic constraints give
Eq.~\ref{eq:cpgsv21_normal_range_reduction_block_word_trace_identity}. Trace
separation yields
Eq.~\ref{eq:cpgsv21_normal_range_reduction_block_window_matrix_identity}, and
the block-injective commutation lemma yields
Eq.~\ref{eq:cpgsv21_normal_range_reduction_one_site_commutation}. The verified
boundary-restriction equality completes the periodic comparison. At
$N=L_0+1$, the cyclic change-of-cut argument proves
Eq.~\ref{eq:cpgsv21_normal_range_reduction_minimal_ring_ground_space}.
Therefore the formal theorem covers the full source range $L_0+1\leq N$.

\bibliographystyle{alpha}
\bibliography{references}

\end{document}
